Write
\[
\mathcal R_{n,d}(\rho)=
\inf_{\widetilde\theta_n}
\sup_{(P_{0,n},P_{1,n})\in\mathcal M_{n,d}(s_y,s_z,\rho)}
\{E(\widetilde\theta_n-\theta_n)^2\}^{1/2}.
\]
We prove the upper and lower inequalities separately.

First we record the relation between the causal endpoint and the radius used to define the statistical class.  For a fixed \(z\), let
\[
F_{a,z}(y)=\int_0^yq_{a,n}(v,z)\,dv,
\qquad Q_{a,z}=F_{a,z}^{-1},
\qquad H_z=F_{1,z}-F_{0,z}.
\]
The inverses exist and are continuous because
\(q_{a,n}\ge c>0\) by
\(\mathrm{ass{:}conditional\text{-}density\text{-}bounds}\).
The increasing coupling \((Q_{0,z}(U),Q_{1,z}(U))\), with \(U\) uniform, minimizes squared cost: indeed, whenever \(x\le x'\) and \(y\le y'\), direct expansion gives
\[
(x-y)^2+(x'-y')^2
\le (x-y')^2+(x'-y)^2.
\]
Successive uncrossing therefore gives the one-dimensional formula
\[
W_2^2\{Q_{0,n}(\cdot\mid z),Q_{1,n}(\cdot\mid z)\}
=\int_0^1\{Q_{1,z}(v)-Q_{0,z}(v)\}^2\,dv.                 \tag{1}
\]
Let \(u_z=\mathcal G_{n,z}\Delta_n(\cdot,z)\).  Since
\(\int_0^1\Delta_n(y,z)\,dy=0\), integration of the defining Neumann equation from \(0\) to \(y\) gives
\[
-\bar q_n(y,z)u_z'(y)=H_z(y).                              \tag{2}
\]
The boundary term vanishes because \(H_z(0)=H_z(1)=0\).  A second integration by parts consequently yields
\[
\int_0^1\Delta_n(y,z)u_z(y)\,dy
=\int_0^1\bar q_n(y,z)\{u_z'(y)\}^2\,dy
=\int_0^1\frac{H_z(y)^2}{\bar q_n(y,z)}\,dy.              \tag{3}
\]
For every \(v\in(0,1)\), the mean-value theorem and
\(F_{a,z}\{Q_{a,z}(v)\}=v\) give a point \(\zeta_{v,z}\) between the two quantiles such that
\[
H_z\{Q_{0,z}(v)\}
=q_{1,n}(\zeta_{v,z},z)\{Q_{0,z}(v)-Q_{1,z}(v)\}.          \tag{4}
\]
Changing variables \(v=F_{0,z}(y)\) in (1), and then using
\(c\le q_{a,n},\bar q_n\le C\) in (3)--(4), proves
\[
\frac{c^2}{C^2}
\int_0^1\Delta_n\mathcal G_{n,z}\Delta_n\,dy
\le W_2^2\{Q_{0,n}(\cdot\mid z),Q_{1,n}(\cdot\mid z)\}
\le
\frac{C^2}{c^2}
\int_0^1\Delta_n\mathcal G_{n,z}\Delta_n\,dy.            \tag{5}
\]
The common-marginal assumption permits integration of (5) against the single measure \(R_n\).  By the definitions of \(\rho_n\) and \(\theta_n\), followed by
\(\mathrm{ass{:}local\text{-}equality\text{-}radius}\),
\[
\frac{c^2}{C^2}\rho_n^2\le\theta_n
\le\frac{C^2}{c^2}\rho_n^2
\le\frac{C^2}{c^2}\rho^2.                                \tag{6}
\]
Thus the identically zero estimator has uniform root-mean-square risk at most \((C^2/c^2)\rho^2\).

Put
\[
g_n=n^{-2s_z/(2s_z+d)}
\varepsilon_{\mathrm{pil},n}^{d/(2s_z+d)},
\qquad
p_n=\varepsilon_{\mathrm{pil},n}^{2+a_y}
+\rho^{2a_y/(2+3/s_y+d/s_z)}\varepsilon_{\mathrm{pil},n}^2.
\]
The exact expansion in
\(\mathrm{thm{:}uniform\text{-}second\text{-}order\text{-}expansion}\)
uses the observable sample map in
\(\mathrm{def{:}anisotropic\text{-}estimator}\) and gives
\(V_{L,n}\lesssim\rho_n^2\), with no reverse-variance premise.  At the resolutions in
\(\mathrm{def{:}anisotropic\text{-}tuning}\), it therefore gives
\[
\sup_{P\in\mathcal M_{n,d}(s_y,s_z,\rho)}
\|\widehat\theta^{\mathrm{an}}_{J_y^*,J_z^*,n}-\theta_n\|_2
\lesssim \rho n^{-1/2}+g_n+p_n.                            \tag{7}
\]
For every fixed \(0<\tau<\min\{s_y,s_z\}\),
\(\mathrm{lem{:}anisotropic\text{-}smoothcot\text{-}median\text{-}envelope}\)
gives
\[
\sup_{P\in\mathcal M_{n,d}(s_y,s_z,\rho)}
\|\widehat\theta^{\mathrm{sm,med}}_{\tau,n}-\theta_n\|_2
\le C_\tau n^{-a_\tau}(\log n)^{a_\tau}.                  \tag{8}
\]
When \(s_z\ge1\),
\(\mathrm{lem{:}direct\text{-}cot\text{-}median\text{-}envelope}\)
similarly gives
\[
\sup_{P\in\mathcal M_{n,d}(s_y,s_z,\rho)}
\|\widehat\theta_n^{\mathrm{dir,med}}-\theta_n\|_2
\lesssim n^{-1/(d+2)}(\log n)^{1+1/(d+2)}.                 \tag{9}
\]
For fixed class indices, choose before observing the sample the estimator whose deterministic guarantee among (6)--(9) is smallest.  This is one data-dependent map with a law-independent choice.  Absorbing its fixed constants into \(C_{1,\tau}\) proves the displayed upper bound for \(\mathcal R_{n,d}(\rho)\).

We next prove the converse.  At the resolution
\(2^{J_z}\asymp n^{2/(4s_z+d)}\),
\(\mathrm{lem{:}explicit\text{-}mixture\text{-}checks}\) places the null and every component of
\(\mathrm{def{:}rademacher\text{-}family}\) in the same class.  By decreasing the fixed \(c_0\), if necessary, that lemma gives constants \(\eta<1\) and \(c_M>0\) such that
\[
\chi^2(\overline{\mathbb P}_n,\mathbb P_n^0)\le\eta,
\qquad
\theta_n(\omega)-\theta_n(0)
\ge c_M\min\{\rho^2,r_{n,d}(s_z)\}.                        \tag{10}
\]
Cauchy--Schwarz applied to the mixture likelihood ratio gives
\[
\|\overline{\mathbb P}_n-\mathbb P_n^0\|_{\mathrm{TV}}
\le\frac12\chi^2(\overline{\mathbb P}_n,\mathbb P_n^0)^{1/2}
\le\frac12\sqrt\eta.                                      \tag{11}
\]
Given any estimator, threshold it halfway between the null endpoint and the lower endpoint separation in (10).  On either testing error its absolute estimation error is at least half that separation.  Averaging the alternative risks over \(\omega\), using (11), and using
\(E|X|\le\{E(X^2)\}^{1/2}\), yields
\[
\mathcal R_{n,d}(\rho)
\ge c_2\min\{\rho^2,r_{n,d}(s_z)\}.                        \tag{12}
\]

It remains to obtain the regular term.  Extend the symmetric cosine family in
\(\mathrm{def{:}no\text{-}covariate\text{-}witness}\) constantly in \(z\), and denote its endpoint by \(\Theta(t)\).  Here
\(\Delta_t(y)=t\cos(2\pi y)\), \(\bar q_t=1\), and direct solution of the Neumann equation gives
\[
\mathcal G\Delta_t(y)=\frac{t}{4\pi^2}\cos(2\pi y),
\qquad
\rho(t)^2=\int_0^1\Delta_t\mathcal G\Delta_t\,dy
=\frac{t^2}{8\pi^2}.                                        \tag{13}
\]
For \(\rho\ge K n^{-1/2}\), choose fixed sufficiently small positive constants and set
\[
t=c_7\min\{\rho,t_0\},
\qquad t'=t+c_8n^{-1/2}.                                    \tag{14}
\]
Equation (13) shows that both laws have radius at most \(\rho\) once \(c_7,c_8\) are small and \(K\) is large.  Positivity, normalization, H\"older membership, and bounded product Hellinger distance follow from
\(\mathrm{ass{:}holder\text{-}radius\text{-}interior}\) and
\(\mathrm{lem{:}explicit\text{-}mixture\text{-}checks}\).

Implicit differentiation of
\(F_{a,t}\{Q_{a,t}(v)\}=v\) gives
\[
\partial_tQ_{a,t}(v)
=-\frac{(2a-1)\sin\{2\pi Q_{a,t}(v)\}}
{4\pi q_{a,t}^{0}\{Q_{a,t}(v)\}}.                         \tag{15}
\]
The denominator is uniformly bounded below.  Thus (15) and dominated differentiation make \(\Theta''\) continuous near zero.  By
\(\mathrm{lem{:}cosine\text{-}second\text{-}variation}\),
\(\Theta'(0)=0\) and \(\Theta''(0)=1/(4\pi^2)\).  Shrinking \(t_0\) if necessary, the mean-value theorem applied first to \(\Theta'\) and then to \(\Theta\) gives
\[
|\Theta(t')-\Theta(t)|
\ge c_3\min\{\rho,t_0\}n^{-1/2}
\ge c_3t_0\rho n^{-1/2},                                  \tag{16}
\]
where the last inequality uses \(0\le\rho\le1\).  The two-point testing argument used above, now with the Hellinger affinity supplied by the mixture-check lemma, proves
\[
\mathcal R_{n,d}(\rho)\ge c_4\rho n^{-1/2}
\quad\text{whenever }\rho\ge K n^{-1/2}.                   \tag{17}
\]

Set \(r=r_{n,d}(s_z)=n^{-b}\), where
\(b=4s_z/(4s_z+d)<1\).  If \(\rho<K n^{-1/2}\), then, uniformly for all sufficiently large \(n\),
\(\rho^2\le K^2n^{-1}=o(r)\), so (12) is of order \(\rho^2\), which equals the target minimum.  If \(\rho\ge K n^{-1/2}\) and \(\rho^2\le r\), (12) again supplies \(\rho^2\).  Finally, if \(\rho\ge K n^{-1/2}\) and \(\rho^2>r\), (12) and (17) imply a lower bound by a constant times
\(\max\{r,\rho n^{-1/2}\}\), which is at least half of
\(r+\rho n^{-1/2}\).  In all three cases,
\[
\mathcal R_{n,d}(\rho)
\gtrsim
\min\{\rho^2,\rho n^{-1/2}+r_{n,d}(s_z)\}
=\mathfrak R_{n,d}(s_y,s_z,\rho).                           \tag{18}
\]
This proves the lower inequality.

For completeness, we verify the claimed strict gap and phase boundaries.  Write
\[
\varepsilon_{\mathrm{pil},n}=(\log n/n)^e,
\qquad
e=\frac{1}{2+s_y^{-1}+ds_z^{-1}}.
\]
Because \(d\ge1\) and \(s_y,s_z>0\),
\[
e<\frac{s_z}{2s_z+d}<\frac{2s_z}{4s_z+d}.                  \tag{19}
\]
The polynomial exponent of \(g_n\) is
\((2s_z+ed)/(2s_z+d)\), and direct subtraction gives
\[
\frac{4s_z}{4s_z+d}-\frac{2s_z+ed}{2s_z+d}
=\frac{d\{2s_z-e(4s_z+d)\}}
{(4s_z+d)(2s_z+d)}>0.                                       \tag{20}
\]
The logarithmic factor in \(g_n\) has positive power, so (20) proves
\[
n^{-4s_z/(4s_z+d)}=o(g_n).                                  \tag{21}
\]
Since \(x\mapsto x/(2x+d+1)\) is strictly increasing and \(\tau<s_z\),
\[
a_\tau<\frac{s_z}{2s_z+d+1}
<\min\left\{\frac{4s_z}{4s_z+d},\frac12\right\}.         \tag{22}
\]
When \(s_z\ge1\), direct cross-multiplication also gives
\[
\frac1{d+2}<\min\left\{\frac{4s_z}{4s_z+d},\frac12\right\}. \tag{23}
\]
Thus the global fallbacks can sharpen the upper envelope but do not establish the intrinsic quadratic exponent.

Finally, solving \(\rho_{Q,n}^2=r_{n,d}(s_z)\) and
\(\rho_{L,n}n^{-1/2}=r_{n,d}(s_z)\) gives
\[
\rho_{Q,n}=n^{-2s_z/(4s_z+d)},
\qquad
\rho_{L,n}=n^{(d-4s_z)/\{2(4s_z+d)\}}.                     \tag{24}
\]
The second boundary tends to zero exactly when \(s_z>d/4\).  If \(s_z<d/4\), the intrinsic quadratic lower scale is slower than \(n^{-1/2}\); at the equality boundary the first derivative vanishes by
\(\mathrm{lem{:}cosine\text{-}second\text{-}variation}\).  At \(s_z=d/4\), that lower scale is \(n^{-1/2}\) but still arises from the degenerate quadratic component rather than a positive first-order variance.  Hence for \(s_z\le d/4\) no uniformly positive-variance root-\(n\) description covers the boundary.  Equations (6)--(24) prove every assertion of the theorem.